% theme: gymnosperms_and_spore_plants
\MajorQuestion{植物のなかま（2）}
\SetRowSpaceQanda

\Instruction{次の問いに答えなさい。}
\QQ{マツの新芽の先に集まり、りん片に胚珠がついている花を何というか。}{雌花}
\QQ{種子をつくらない植物の多くは、どのような場所に生息しているか。}{日当たりが悪く、湿り気の多いところ}
\QQ{マツのように、胚珠がむき出しになっている植物を\NamedBlank{gymnosperm}という。}{裸子植物}
\QQ{マツの花粉や種子は、おもに何によって運ばれるか。}{風}
\QQ{マツの新芽のつけ根に集まり、りん片で花粉をつくる花を\NamedBlank{malecone}という。}{雄花}

\Instruction{\strut}
\QQ{種子をつくらない植物が、なかまをふやすためにつくる小さな粒を\NamedBlank{spore}という。}{胞子}
\QQ{根・茎・葉の区別があり、葉の裏に袋をつけることが多い植物を\NamedBlank{fern}という。}{シダ植物}
\QQ{花を咲かせ、種子をつくってふえる植物を何というか。}{種子植物}
\QQ{\RefBlank{malecone}のりん片にあり、花粉をつくる袋を何というか。}{花粉のう}
\QQ{根・茎・葉の区別がなく、雌株と雄株がある植物を\NamedBlank{moss}という。}{コケ植物}

\Instruction{\strut}
\QQ{胚珠が子房の中にある植物を何というか。}{被子植物}
\QQ{\RefBlank{fern}には、根・茎・葉の区別があるか、ないか。}{ある}
\QQ{\RefBlank{moss}の例を1つ答えなさい。}{ゼニゴケ（スギゴケ、コスギゴケ）}
\QQ{\RefBlank{gymnosperm}には、めしべの根元にある何が存在しないか。}{子房}
\QQ{地中を横にのびる、\RefBlank{fern}の茎を何というか。}{地下茎}

\Instruction{\strut}
\QQ{\RefBlank{moss}は体のどの部分から水を吸収するか。}{体の表面全体}
\QQ{\RefBlank{fern}の葉の裏にできる、\RefBlank{spore}をつくる袋を\NamedBlank{sporangium}という。}{胞子のう}
\QQ{\RefBlank{fern}の例を、イヌワラビ以外に1つ答えなさい。}{ゼンマイ（スギナ、ノキシノブ）}
\QQ{\RefBlank{moss}にある、根のように見える部分を何というか。}{仮根}
\QQ{\RefBlank{moss}で、\RefBlank{sporangium}ができる株を何というか。}{雌株}
